En el Capítulo 1, Antecedentes, se presenta el planteamiento del problema, la propuesta de solución, los objetivos generales y específicos, así como la justificación del proyecto. Se describe la problemática relacionada con la selección de alineaciones y sustituciones, además de delimitar el alcance del modelo propuesto.

En el Capítulo 2, Estado del arte, se revisan investigaciones relacionadas con el análisis del rendimiento deportivo, las interacciones entre jugadoras, la evaluación de alineaciones y la aplicación de métodos matemáticos y computacionales en el baloncesto. 

En el Capítulo 3, Marco teórico, se describen los conceptos necesarios para comprender el desarrollo del proyecto. Se incluyen fundamentos del baloncesto profesional femenino, estadísticas individuales y colectivas, métricas avanzadas de rendimiento, métodos de aprendizaje automático, modelos estadísticos, técnicas de optimización matemática y medidas para la evaluación de modelos predictivos.

En el Capítulo 4, Análisis, se presenta la metodología en cascada adoptada para organizar las actividades del proyecto. También se describen las herramientas y bibliotecas consideradas, las fuentes de información, la estructura de los conjuntos de datos y las variables necesarias para las etapas posteriores. Finalmente, se incluyen los análisis de factibilidad y riesgos asociados al desarrollo del modelo.

En el Capítulo 5, Preprocesamiento, se describen las actividades relacionadas con la recopilación, exploración, limpieza, transformación e integración de los datos históricos de la WNBA. Se consideran los procedimientos necesarios para verificar la consistencia de los registros, identificar las alineaciones participantes y preparar la información que servirá como base para el análisis de jugadoras y quintetos.

En el Capítulo 6, Desarrollo y evaluación del modelo, se presenta el diseño preliminar de la solución propuesta. Se describe un procedimiento dividido en dos fases: la caracterización de jugadoras y el análisis de sus interacciones, seguido de la estimación del rendimiento y la selección de alineaciones frente a un quinteto rival. Se establecen las consideraciones iniciales para la formulación matemática, la comparación de alternativas y la posterior evaluación del modelo.